\documentclass[10pt]{article}

% Specify the margins and text size.
\setlength{\textwidth}{6.4in}
\setlength{\textheight}{9.5in}
\setlength{\oddsidemargin}{0pt}
\setlength{\evensidemargin}{0pt}
\setlength{\topmargin}{0pt}
\setlength{\hoffset}{.05in}
\setlength{\voffset}{-.8in}

\setlength{\parskip}{5pt}
\setlength{\parindent}{0pt}

% Load some fonts and symbol packages
\usepackage{latexsym}
\usepackage{pifont}       % contains 'star' symbol for counterinsurgency handbook title
\usepackage{yfonts} 
\usepackage{amsmath}
\usepackage{amsfonts}

\usepackage{graphicx}     % actually, this is loaded by pstricks
\usepackage[T1]{fontenc}
\usepackage{ifthen}
\usepackage{pstricks,pst-grad,pst-text,pst-node,multido,pst-plot,calc,pst-3dplot}
%\usepackage[all]{xy}
%\usepackage{animate}

% The hyperref package inserts links into the pdf.
\definecolor{MyLinkColor}{rgb}{.1,.2,1}
\definecolor{MyCiteColor}{rgb}{.1,1,.2}
\definecolor{MyURLColor}{rgb}{.4,.4,.4}
\usepackage[backref=true,pagebackref=false,hyperindex,colorlinks=true,
  linkcolor=MyLinkColor,urlcolor=MyURLColor]{hyperref}

\definecolor{Red}{rgb}{1,0,0}
\definecolor{myblue}{rgb}{0,0,0.6}
\definecolor{Blue}{rgb}{0,0,1}
\definecolor{Green}{rgb}{0,1,0}
\definecolor{magenta}{rgb}{1,0,.6}
\definecolor{lightblue}{rgb}{0,.5,1}
\definecolor{lightpurple}{rgb}{.6,.4,1}
\definecolor{gold}{rgb}{.6,.5,0}
\definecolor{orange}{rgb}{1,0.4,0}
\definecolor{hotpink}{rgb}{1,0,0.5}
\definecolor{newcolor2}{rgb}{.5,.3,.5}
\definecolor{newcolor}{rgb}{0,.3,1}
\definecolor{newcolor3}{rgb}{1,0,.35}
\definecolor{darkgreen1}{rgb}{0, .35, 0}
\definecolor{darkgreen}{rgb}{0, .6, 0}
\definecolor{darkred}{rgb}{.75,0,0}
\definecolor{grayA}{rgb}{0.95,0.95,0.95}
\definecolor{grayB}{rgb}{0.98,0.98,0.98}

\xdefinecolor{olive}{cmyk}{0.64,0,0.95,0.4}
\xdefinecolor{purpleish}{cmyk}{0.75,0.75,0,0}


% The tweaklist package is something I found on the web.  It provides a simple interface
% for making changes to spacing used in the itemize and enumerate environments.  Comment
% this out if you don't care to use tweaklist.
\usepackage{tweaklist}
\renewcommand{\itemhook}{\setlength{\parskip}{2pt}\setlength{\parsep}%
{1pt}\setlength{\topsep}{0pt}\setlength{\itemsep}{0pt}}

\newcommand{\U}{\underline{\hspace{5pt}}}

\usepackage{listings}
\newcounter{EX}\setcounter{EX}{1}
\newcommand{\EXERCISE}{\arabic{EX}.\stepcounter{EX} }

\begin{document}
\pagestyle{empty}
\lstset{language=R, showspaces=false, showstringspaces=false}

\href{http://www.su.edu}{\includegraphics[height=1.75cm]{sulogo.eps}}
\vspace{-1.69cm}

{\small \hfill
\begin{tabular}{cl}
& Math 207\\ & Introduction to Statistics\\
\end{tabular}
}
\setlength{\baselineskip}{1.05\baselineskip}

\begin{center}
\textbf{\large  Controlled Experiment Exercises}
\end{center}

1. The Salk vaccine field trial tables are shown below.\vspace{-3pt}
\begin{center}{%\footnotesize
\begin{tabular}{cc}
Randomized controlled          & \\
double-blind experiment        & NFIP study\\
   \begin{tabular}{lcc}\hline
                 & Size      & Rate${}^*$\vphantom{\LARGE Y}\\
      \color{darkgreen}Treatment  & \color{darkgreen}200,000   & {\color{blue}\textbf{28}}\\
      \color{darkgreen}Control    & \color{darkgreen}200,000   & {\color{blue}\textbf{71}}\\
      No consent & 350,000   & 46\\
   \end{tabular}
&
   \begin{tabular}{lcc}
\hline
                 & Size        &  Rate${}^*$\vphantom{\LARGE Y}\\
      \color{darkgreen}Grade 2 (vaccine)\hfill        & \color{darkgreen}225,000 & 25\\
      \color{gray}Grades 1 and 3 (control)        & \color{gray}725,000 & {\color{red}\textbf{54}}\\
      Grade 2 (no consent)\hfill     & 125,000 & {\color{red}\textbf{44}}\\
   \end{tabular}
\end{tabular}}
\end{center}\vspace{-2pt}

\hspace{10pt} a) In the randomized controlled experiment, which numbers indicate the effectiveness of the vaccine?
\vspace{.2in}

\hspace{10pt} b) In the NFIP study, which numbers indicate the effectiveness of the vaccine?
\vspace{.2in}

\hspace{10pt} c) In the randomized controlled experiment, neither the no-consent group nor the control group 
received the vaccine.  Why did their polio rates differ so significantly?
\vspace{.5in}

\hspace{10pt} d) In the NFIP study, the control and no-consent groups contracted polio at different rates?  Why?
\vspace{.5in}

\hspace{10pt} e) The sizes of the groups differ.  How does this impact our ability to interpret the results?
\vspace{.5in}

2. Data from the Salk vaccine field trial suggest that in 1954, the school districts in the NFIP 
trial and in the randomized controlled experiment had similar 
exposures to the polio virus.

\hspace{10pt} a) The data also show that children in the two vaccine groups (for the randomized
controlled experiment and the NFIP design) came from families with similar incomes and educational
backgrounds.  Which two numbers confirm this finding?
\vspace{.3in}

\hspace{10pt} b) The data show that children in the two no-consent groups had similar family backgrounds.
Which pair of numbers confirm this finding?
\vspace{.3in}

\hspace{10pt} c) The data show that children in the two control groups had different family
backgrounds.  Which pair of numbers in the table confirm this finding?
\vspace{.3in}

\hspace{10pt} d) In the NFIP study, neither the control group nor the no-consent group got the 
vaccine.  yet the no-consent group had a lower rate of polio.  Why?
\vspace{.5in}

\hspace{10pt} e) To show that the vaccine works, someone wants to compare the 44/100,000 in the
NFIP study with the 25/100,000 in the vaccine group.  What's wrong with this idea?
\vspace{.6in}

3. The Salk vaccine field trials were conducted only in certain experimental areas (school districts),
selected by the Public Health Service in consultation with local officials.  
In these areas, there were about 3 million children in grades 1, 2, or 3; and there were about 
11 million children in those grades in the United States.  In the experimental areas,
the incidence of polio was about 25\% higher than in the rest of the country.  Did the Salk
vaccine field trials cause children to get polio instead of preventing it?  
\vspace{1in}

4. In the randomized controlled Salk vaccine trial, those children whose parents refused to 
participate got polio at the rate of 46 per 100,000.  On the other hand, those children whose
parents consented to participation got polio at the slightly higher rate of 49 per 100,000
in the treatment group and control group taken together.  
Suppose that this field trial was repeated the following year.  On the basis of the figures, some parents
refused to allow their children to participate in the experiment and be exposed to this higher riskof polio.
Were they right?
\vspace{1in}

5. In and educational study, students choose to enroll in either an online statistics course or a face-to-face
course.  At the end of the semester, all students are given a common final exam. Student grades on the final 
exam are used to compare the effectiveness of the course formats.

\hspace{10pt} a) Is this a controlled experiment or an observational study?
\vspace{.3in}


\hspace{10pt} b) What are some characteristics, other than course format, that differ between the two 
groups?
\vspace{.5in}

6. A medical device company is testing a system that can turn hospital patients in their beds at regular intervals
to avoid bed sores.  500 patients agree to participate in the study.  Of these 250 are randomly assigned
to the treatment group and will be turned by the device.  The other 250 will be turned manually by hospital 
staff.  Health care providers working with the treatment group, however, do not use the device consistently.
Some use it to turn patients every two hours as directed while others sometimes use the device but other times
turn patients manually.  

\hspace{10pt} a) Is this a controlled experiment or an observational study?
\vspace{.3in}

\hspace{10pt} b) Suggest some reasons why providers may have used the device with different frequencies.
\vfill
\eject

\end{document}
